\chapter{总结与展望}
\label{chap:conclusion}

\section{工作总结}

本实验以 Oxford-IIIT Pet Dataset 为研究对象，系统性地探索了图像语义分割任务从 baseline 构建到多层次优化的完整流程。

在 baseline 阶段，实现了完整的 U-Net 分割 pipeline，涵盖数据加载、模型训练和评估，在测试集上达到 mIoU 0.71，建立了可复现的性能基准。在此基础上，从数据、模型、训练策略和推理四个层次进行了系统化优化：数据侧引入 Albumentations 高级增强和类别权重 $w=[1,1,3]$；模型侧使用 ImageNet 预训练的 ResNet34 替换手写编码器，参数量从 31M 降至 21M；训练策略方面采用 Focal Tversky Loss、AdamW 优化器、CosineAnnealingWarmRestarts 学习率调度和 Early Stopping；推理侧通过 TTA 水平翻转平均进一步提升性能。最终模型 mIoU 达到 0.81，累计提升 10.5\%，轮廓类 IoU 从 0.47 提升至 0.61。

\section{关键发现}

预训练 Encoder 是性能提升的最大贡献者，独立贡献 4.0\% 的 mIoU 提升。ResNet34 的低层特征——边缘检测器和纹理识别器——可直接迁移到宠物分割任务，使模型无需从零学习基础视觉特征，同时参数量反而减少了 32\%。这一结果充分验证了 ImageNet 预训练特征在下游任务中的泛化能力。

在类别不平衡的应对上，类别权重与 Focal Tversky Loss 的组合使轮廓类 IoU 提升了 14.3\%。$\alpha=0.7 > \beta=0.3$ 的设置使模型更关注漏检问题，这对仅占 3.2\% 像素的轮廓类尤为关键。消融实验还揭示了各优化项之间的协同效应：累积效果 10.5\% 大于独立贡献之和，说明预训练 Encoder 提供的更好特征表示能够放大数据增强和损失函数优化的效果。

\section{局限性与未来工作}

当前模型仍存在若干局限。在强遮挡场景下，模型无法推断被遮挡部分的轮廓，导致前景区域不完整。当宠物毛色与背景颜色接近时，边界模糊导致误分类。极端光照条件下纹理特征丢失，分割精度明显下降。此外，轮廓类 IoU 虽然提升显著，但 0.61 的绝对值仍远低于前景和背景，2--3 像素宽的边界本身就极难精确分割。TTA 使推理时间约增加一倍，不适合实时应用场景。

未来可从以下方向继续改进。模型架构方面，可尝试 EfficientNet-B4 或 Swin Transformer 等更强的预训练 Backbone，或引入 FPN、DeepLabV3+ 等多尺度架构增强对不同尺度目标的感知。在解码器中加入 CBAM、SE 等注意力模块有望进一步增强边缘感知能力。训练范式方面，半监督学习可利用未标注数据进一步提升性能。部署方面，模型剪枝和知识蒸馏可实现移动端的轻量化部署。本实验的方法论具有较好的可迁移性，可推广至医学图像分割、自动驾驶场景理解和遥感图像分析等领域。

\section{个人收获}

这次实验是我第一次完整地经历一个语义分割项目从零到一的全过程，收获远超预期。

在技术层面，我对 U-Net 的编码器-解码器架构有了更深的理解。最初手写 U-Net 时，我对跳跃连接的作用只停留在理论认知，直到观察到 baseline 的预测结果中边缘模糊的问题，才真正体会到多尺度特征融合对空间细节恢复的重要性。引入预训练 ResNet34 后，模型性能的跳跃式提升让我深刻认识到迁移学习的威力——在数据量有限的场景下，利用大规模数据集上学到的通用视觉特征，远比从零训练更加高效。

在实验方法论上，消融实验的设计让我学会了如何科学地量化每个优化项的贡献。起初我倾向于一次性叠加所有优化策略，但这样无法判断哪些改进真正有效。通过逐项添加和独立对比，我不仅明确了各策略的贡献排序，还发现了它们之间的协同效应，这种系统化的实验思维对未来的研究工作非常有价值。

在工程实践上，我积累了不少经验。Albumentations 库的图像-mask 同步变换机制让我意识到数据增强在分割任务中需要特别注意标注对齐问题。调试 Focal Tversky Loss 的超参数时，我体会到损失函数设计并非简单的公式替换，而是需要结合具体任务的特点进行针对性调整。学习率 warmup 策略的引入也让我理解了预训练权重的脆弱性——过大的初始梯度会破坏已有的特征表示，这一点在实验中得到了验证。

回顾整个实验过程，从问题分析到假设提出，从实验验证到结果解释，我逐步建立起了较为完整的科研思维框架。这种从数据出发、用实验说话的方法论，将成为我未来学习和研究中的重要基础。

\section{致谢}

感谢 Oxford Visual Geometry Group 提供的 Pet Dataset，以及开源社区维护的 PyTorch、Segmentation Models PyTorch、Albumentations 等工具库。
